

%We benchmark Equiformer on two quantum properties prediction datasets, QM9~\cite{qm9_2, qm9_1} and OC20~\cite{oc20}.
%Our implementation is based on PyTorch~\cite{pytorch}, PyG~\cite{pytorch_geometric}, \texttt{e3nn}~\cite{e3nn} and \texttt{timm}~\cite{timm}.


\begin{table}[t]
\begin{adjustwidth}{0cm}{0cm}
\centering
\resizebox{\ifdim\width>\textwidth\textwidth\else\width\fi}{!}{
\begin{tabular}{llcccccccccccc}
%\cline{7-11}
\toprule[1.2pt]
& Task & $\alpha$ & $\Delta \varepsilon$ & $\varepsilon_{\text{HOMO}}$ & $\varepsilon_{\text{LUMO}}$ & $\mu$ & $C_{\nu}$ & $G$ & $H$ & $R^2$ & $U$ & $U_0$ & ZPVE \\ 
%Methods & Units & bohr$^3$ & meV & meV & meV & D & cal/mol K & meV & meV & bohr$^3$ & meV & meV & meV\\
Methods & Units & $a_0^3$ & meV & meV & meV & D & cal/mol K & meV & meV & $a_0^2$ & meV & meV & meV\\
\midrule[1.2pt]

NMP~\citep{neural_message_passing_quantum_chemistry}$^{\dagger}$ & & .092 & 69 & 43 & 38 & .030 & .040 & 19 & 17 & .180 & 20 & 20 & 1.50 \\
SchNet~\citep{schnet} & & .235 & 63 & 41 & 34 & .033 & .033 & 14 & 14 & .073 & 19 & 14 & 1.70 \\
Cormorant~\citep{cormorant}$^{\dagger}$ & & .085 & 61 & 34 & 38 & .038 & .026 & 20 & 21 & .961 & 21 & 22 & 2.03 \\
LieConv~\citep{lieconv}$^{\dagger}$ & & .084 & 49 & 30 & 25 & .032 & .038 & 22 & 24 & .800 & 19 & 19 & 2.28 \\
DimeNet++~\citep{dimenet_pp} & & \textbf{.044} & 33 & 25 & 20 & .030 & .023 & 8 & 7 & .331 & 6 & 6 & 1.21 \\
TFN~\citep{tfn}$^{\dagger}$ & & .223 & 58 & 40 & 38 & .064 & .101 & - & - & - & - & - & - \\
\multicolumn{2}{l}{SE(3)-Transformer~\citep{se3_transformer}$^{\dagger}$} & .142 & 53 & 35 & 33 & .051 & .054 & - & - & - & - & - & - \\
EGNN~\citep{egnn}$^{\dagger}$ & & .071 & 48 & 29 & 25 & .029 & .031 & 12 & 12 & .106 & 12 & 11 & 1.55 \\

PaiNN~\citep{painn} & & .045 & 46 & 28 & 20 & .012 & .024 & \textbf{7.35} & \textbf{5.98} & .066 & \textbf{5.83} & \textbf{5.85} & 1.28 \\

TorchMD-NET~\citep{torchmd_net} & & .059 & 36 & 20 & 18 & \textbf{.011} & .026 & 7.62 & 6.16 & \textbf{.033} & 6.38 & 6.15 & 1.84 \\

SphereNet~\citep{spherenet} & & .046 & 32 & 23 & 18 & .026 & \textbf{.021} & 8 & 6 & .292 & 7 & 6 & \textbf{1.12} \\

SEGNN~\citep{segnn}$^{\dagger}$ & & .060 & 42 & 24 & 21 & .023 & .031 & 15 & 16 & .660 & 13 & 15 & 1.62 \\

%Equiformer& 0.5093& 0.6285 & 0.5053 & 0.5556 & 0.5497 & 5.04 & 2.85 & 4.90 & 3.04 & \\
EQGAT~\citep{eqgat} & & .053 & 32 & 20 & 16 & \textbf{.011} & .024 & 23 & 24 & .382 & 25 & 25 & 2.00 \\
\midrule[0.6pt]

Equiformer & & .046 & \textbf{30} & \textbf{15} & \textbf{14} & \textbf{.011} & .023 & 7.63 & 6.63 & .251 & 6.74 & 6.59 & 1.26 \\
%\midrule[0.6pt]
%Equiformer + NAT & 0.4737 & 0.7245 & 0.4862 & 0.6493 & 0.5834 & 6.01 & 2.45 & 5.41 & 2.71 & \\
\bottomrule[1.2pt]

\end{tabular}
}
\vspace{-3mm}
\end{adjustwidth}
\caption{\textbf{MAE results on QM9 testing set.}  
%\note{Dropout rate = 0.1.}
%%%$\dagger$ denotes using different training, validation, testing data partitions. % as mentioned in SEGNN~\cite{segnn}.
$\dagger$ denotes using different data partitions.
%$\ddagger$ denotes results from SE(3)-Transformer~\cite{se3_transformer}.
}
\vspace{-3mm}
\label{tab:qm9_result}
\end{table}


\vspace{-2.5mm}
\section{Experiment}
\vspace{-3mm}

%%%Our implementation is based on PyTorch~\citep{pytorch} (Modified BSD license), PyG~\citep{pytorch_geometric} (MIT license), \texttt{e3nn}~\citep{e3nn} (MIT license), \texttt{timm}~\citep{timm} (Apache-2.0 license), and \texttt{ocp}~\citep{oc20} (MIT license).
We benchmark Equiformer on QM9 (Sec.~\ref{subsec:qm9_exp}), MD17 (Sec.~\ref{subsec:md17}) and OC20 (Sec.~\ref{subsec:oc20}) datasets.
%Moreover, we conduct ablation studies (Sec.~\ref{subsec:ablation_study}) to demonstrate that the proposed equivariant graph attention, consisting of MLP attention and non-linear message passing, improves upon typical dot product attention in Transformer as well as dot product attention in other equivariant Transformers.
\revision{Moreover, ablation studies (Sec.~\ref{subsec:ablation_study}) are conducted to demonstrate that Equiformer with dot prodcut attention and linear message passing has already achieved strong empirical results on QM9 and OC20 datasets and verify that the proposed equivariant graph attention improves upon typical dot product attention in Transformer as well as dot product attention in other equivariant Transformers.}
Additional results of including inversion can be found in Sec.~\ref{appendix:subsec:e3_qm9} and Sec.~\ref{appendix:subsec:e3_oc20}.


\vspace{-3mm}
\subsection{QM9}
\vspace{-3mm}
\label{subsec:qm9_exp}

\paragraph{Dataset.}
%%%The QM9~\citep{qm9_2, qm9_1} dataset (CC BY-NC SA 4.0 license) consisting of 134k small molecules, and the goal is to predict their quantum properties such as energy.
The QM9 dataset~\citep{qm9_2, qm9_1}  (CC BY-NC SA 4.0 license) consists of 134k small molecules, and the goal is to predict their quantum properties.
%The average number of atoms per molecule is 18, and each atom in a molecule comes from the set of H, C, N, O, F atoms and has 3D coordinates.
%%%We follow the data partition used by Cormorant~\cite{cormorant}, which has 100k, 18k and 13k molecules in training, validation and testing sets.
The data partition we use has 110k, 10k, and 11k molecules in training, validation and testing sets.
We minimize mean absolute error (MAE) between prediction and normalized ground truth. % and report MAE on the testing set.


\vspace{-3.25mm}
\paragraph{Training Details.}
Please refer to Sec.~\ref{appendix:subsec:qm9_training_details} in appendix for details on architecture, hyper-parameters and training time.


\vspace{-3.25mm}
\paragraph{Results.}
%%%We mainly compare with methods trained with the same data partition and summarize the results in Table~\ref{tab:qm9_result}.
We summarize the comparison to previous models in Table~\ref{tab:qm9_result}.
Equiformer achieves overall better results across 12 regression tasks compared to each individual model.
%%%Equiformer achieves the best results on 11 out of 12 tasks among models trained with same data partition.
The comparison to SEGNN, which uses irreps features as Equiformer, demonstrates the effectiveness of combining non-lienar message passing with MLP attention.
Additionally, Equiformer achieves better results for most tasks when compared to other equivariant Transformers, which are SE(3)-Transformer, TorchMD-NET and EQGAT.
This demonstrates a better adaption of Transformers to 3D graphs and the effectiveness of the proposed equivariant graph attention.
We note that for the tasks of $\mu$ and $R^2$, PaiNN and TorchMD-NET use different architectures, which take into account the property of the task.
%\revision{This makes the comparison less clear since we use the same architecture for all tasks.}
In contrast, we use the same architecture for all tasks.
We compare training time in Sec.~\ref{appendix:subsec:qm9_training_time_comparison}.

%%%Besides, the different data partition as denoted by $\dagger$ in Table~\ref{tab:qm9_result} has $10\%$ more molecules in the training set and less data in the testing set, and this can benefit some tasks that are more dependent on data partitions.
%\todo{If there is any justification, can write something like which tasks are more important.}
\begin{table}[t]
\begin{adjustwidth}{0cm}{0cm}
\centering
\resizebox{\ifdim\width>\textwidth\textwidth\else\width\fi}{!}{
\begin{tabular}{lcccccccccccccccc}
\toprule[1.2pt]
                        & \multicolumn{2}{c}{Aspirin} & \multicolumn{2}{c}{Benzene} & \multicolumn{2}{c}{Ethanol} & \multicolumn{2}{c}{Malonaldehyde} & \multicolumn{2}{c}{Naphthalene} & \multicolumn{2}{c}{Salicylic acid} & \multicolumn{2}{c}{Toluene} & \multicolumn{2}{c}{Uracil} \\
\cmidrule[0.6pt]{2-17}
Methods                                               & energy       & forces       & energy       & forces       & energy       & forces       & energy          & forces          & energy         & forces         & energy           & forces          & energy       & forces       & energy       & forces      \\
\midrule[1.2pt]
SchNet~\citep{schnet}                          & 16.0         & 58.5         & 3.5          & 13.4         & 3.5          & 16.9         & 5.6             & 28.6            & 6.9            & 25.2           & 8.7              & 36.9            & 5.2          & 24.7         & 6.1          & 24.3        \\
DimeNet~\citep{dimenet}                          & 8.8          & 21.6         & 3.4          & 8.1          & 2.8          & 10.0         & 4.5             & 16.6            & 5.3            & 9.3            & 5.8              & 16.2            & 4.4          & 9.4          & 5.0          & 13.1        \\
PaiNN~\citep{painn}                                                 & 6.9          & 14.7         & -            & -            & 2.7          & 9.7          & 3.9             & 13.8            & 5.0            & 3.3            & 4.9              & 8.5             & 4.1          & 4.1          & 4.5          & 6.0         \\
TorchMD-NET~\citep{torchmd_net}                                           & \textbf{5.3}          & 11.0         & 2.5          & 8.5          & 2.3          & 4.7          & 3.3             & 7.3             & \textbf{3.7}            & 2.6            & \textbf{4.0}              & 5.6             & \textbf{3.2}          & 2.9          & \textbf{4.1}          & 4.1         \\
NequIP ($L_{max} = 3$)~\citep{nequip}                              & 5.7          & 8.0          & -            & -            & \textbf{2.2}          & 3.1          & 3.3             & 5.6             & 4.9            & \textbf{1.7}            & 4.6              & \textbf{3.9}             & 4.0          & \textbf{2.0}          & 4.5          & \textbf{3.3}         \\

\midrule[0.6pt]

Equiformer ($L_{max} = 2$)                                           & \textbf{5.3}          & 7.2          & \textbf{2.2}          & \textbf{6.6}          & \textbf{2.2}          & 3.1          & 3.3             & 5.8             & \textbf{3.7}            & 2.1           & 4.5              & 4.1             & 3.8          & 2.1          & 4.3          & \textbf{3.3}        \\
Equiformer ($L_{max} = 3$)                                           & \textbf{5.3}          & \textbf{6.6}          & 2.5          & 8.1          & \textbf{2.2}          & \textbf{2.9}          & \textbf{3.2}             & \textbf{5.4}             & 4.4            & 2.0           & 4.3              & \textbf{3.9}            & 3.7          & 2.1         & 4.3          & 3.4        \\
\bottomrule[1.2pt]
\end{tabular}
}
\end{adjustwidth}
\vspace{-3mm}
\caption{\textbf{MAE results on MD17 testing set.} Energy and force are in units of meV and meV/$\angstrom$.}
\vspace{-5.5mm}
\label{tab:md17_result}
\end{table}

\vspace{-3mm}
\subsection{MD17}
\vspace{-3mm}
\label{subsec:md17}


\paragraph{Dataset.}
The MD17 dataset~\citep{md17_1, md17_2, md17_3} (CC BY-NC) consists of molecular dynamics simulations of small organic molecules, and the goal is to predict their energy and forces.
We use $950$ and $50$ different configurations for training and validation sets and the rest for the testing set.
Forces are derived as the negative gradient of energy with respect to atomic positions. 
We minimize MAE between prediction and normalized ground truth.

\vspace{-3.25mm}
\paragraph{Training Details.}
Please refer to Sec.~\ref{appendix:subsec:md17_training_details} in appendix for details on architecture, hyper-parameters and training time.


\vspace{-3.25mm}
\paragraph{Results.}
We train Equiformer with $L_{max} = 2 \text{ and } 3$ and summarize the results in Table~\ref{tab:md17_result}.
Equiformer achieves overall better results across 8 molecules compared to each individual model.
%%%\revision{Compared to TorchMD-NET, which is also an equivariant Transformer, Equiformer obtains better MAE results.} 
%since the proposed equivariant graph attention is more expressive and can support vectors of higher degree $L$ (i.e., we use $L_{max} = 2$) instead of restricting to type-$0$ and type-$1$ vectors (i.e., $L_{max} = 1$).}
%%%\revision{The difference lies in the proposed equivariant graph attention, which is more expressive and can support vectors of higher degree $L$ (i.e., we use $L_{max} = 2 \text{ and } 3$) instead of restricting to type-$0$ and type-$1$ vectors (i.e., $L_{max} = 1$).}
Compared to TorchMD-NET, which is also an equivariant Transformer, the difference lies in the proposed equivariant graph attention, which is more expressive and can support vectors of higher degree $L$ (i.e., we use $L_{max} = 2 \text{ and } 3$) instead of restricting to type-$0$ and type-$1$ vectors (i.e., $L_{max} = 1$).
For the last three molecules, although Equiformer with $L_{max} = 2$ achieves lower force MAE but higher energy MAE, we can adjust the weights of energy loss and force loss so that Equiformer achieves lower MAE for both energy and forces as shown in Table~\ref{appendix:tab:md17_torchmdnet} in appendix.
%%%\revision{Compared to NequIP, which also uses irreps features and $L_{max} = 3$, Equiformer with $L_{max} = 2$ achieves overall lower MAE although including higher $L_{max}$ can improve performance and Equiformer uses smaller $L_{max}$.}
Compared to NequIP, which also uses irreps features and $L_{max} = 3$, Equiformer with $L_{max} = 2$ achieves overall lower MAE although including higher $L_{max}$ can improve performance.
When using $L_{max} = 3$, Equiformer achieves lower MAE results for most molecules.
This suggests that the proposed attention can improve upon linear messages even when the size of training sets becomes small.
We compare the training time of NequIP and Equiformer in Sec.~\ref{appendix:subsec:md17_training_time_comparison}.
Additionally, for Equiformer, increasing $L_{max}$ from $2$ to $3$ improves MAE for most molecules except benzene, which results from overfitting.

%\todo{[Training time of NequIP]}
%\todo{[Equiformer with $L_{max} = 3$]}




%\begin{table}[t]
%\centering
%\scalebox{0.7}{
%\begin{tabular}{lccccc|ccccc}
%\cline{7-11}
%\toprule[1.2pt]
% & \multicolumn{5}{c|}{Energy MAE (eV) $\downarrow$} & \multicolumn{5}{c}{EwT (\%) $\uparrow$} \\ 
%\cmidrule[0.6pt]{2-11}
%Methods & ID & OOD Ads & OOD Cat & OOD Both & Average & ID & OOD Ads & OOD Cat & OOD Both & Average \\
%\midrule[1.2pt]
%GNS~\cite{noisy_nodes} & 0.54 & 0.65 & 0.55 & 0.59 & 0.5825 & - & - & - & - & -\\
%GNS + Noisy Nodes~\cite{noisy_nodes} & 0.48 & 0.53 & 0.49 & 0.48 & 0.4950 & - & - & - & - & - \\

%Noisy Nodes~\cite{noisy_nodes} & 0.47 & 0.51 & 0.48 & 0.46 & 0.4800 & - & - & - & - & - \\
%Graphormer~\cite{graphormer_3d} & 0.4329 & 0.5850 & 0.4441 & 0.5299 & 0.4980 & - & - & - & - & -\\
%\midrule[0.6pt]
%Equiformer& 0.5093& 0.6285 & 0.5053 & 0.5556 & 0.5497 & 5.04 & 2.85 & 4.90 & 3.04 & \\
%Equiformer & 0.4222 & 0.5420 & 0.4231 & 0.4754 & 0.4657 & 7.23 & 3.77 & 7.13 & 4.10 & 5.56 \\
%\midrule[0.6pt]
%Equiformer + NAT & 0.4685 & 0.6071 & 0.4746 & 0.5345 & 0.5212 & 5.83 & 3.18 & 5.88 & 3.43 & \\
%\bottomrule[1.2pt]
%\end{tabular}
%}
%\vspace{2mm}
%\caption{\textbf{Results on OC20 IS2RE %\underline{validation} set when IS2RS node-level auxiliary task is adopted during training.} 
%``GNS'' denotes the 50-layer GNS trained without Noisy Nodes data augmentation, and ``Noisy Nodes'' denotes the 100-layer GNS trained with Noisy Nodes.}
%\vspace{-7mm}
%\label{tab:is2re_is2rs_validation}
%\end{table}



\begin{table}[t!]
\begin{adjustwidth}{0cm}{0cm}
\centering
\resizebox{\ifdim\width>\textwidth\textwidth\else\width\fi}{!}{
\begin{tabular}{lccccc|ccccc}
%\cline{7-11}
\toprule[1.2pt]
 & \multicolumn{5}{c|}{Energy MAE (eV) $\downarrow$} & \multicolumn{5}{c}{EwT (\%) $\uparrow$} \\ 
\cmidrule[0.6pt]{2-11}
Methods & ID & OOD Ads & OOD Cat & OOD Both & Average & ID & OOD Ads & OOD Cat & OOD Both & Average \\
\midrule[1.2pt]
CGCNN~\citep{cgcnn} & 0.6149 & 0.9155 & 0.6219 & 0.8511 & 0.7509 & 3.40 & 1.93 & 3.10 & 2.00 & 2.61 \\
SchNet~\citep{schnet} & 0.6387 & 0.7342 & 0.6616 & 0.7037 & 0.6846 & 2.96 & 2.33 & 2.94 & 2.21 & 2.61 \\
DimeNet++~\citep{dimenet_pp} & 0.5621 & 0.7252 & 0.5756 & 0.6613 & 0.6311 & 4.25 & 2.07 & 4.10 & 2.41 & 3.21 \\
PaiNN~\citep{painn} &  0.575 & 0.783 & 0.604 & 0.743 & 0.6763 & 3.46 & 1.97 & 3.46 & 2.28 & 2.79 \\
SpinConv~\citep{spinconv} & 0.5583 & 0.7230 & 0.5687 & 0.6738 & 0.6310 & 4.08 & 2.26 & 3.82 & 2.33 & 3.12 \\
SphereNet~\citep{spherenet} & 0.5625 & 0.7033 & 0.5708 & 0.6378 & 0.6186 & 4.47 & 2.29 & 4.09 & 2.41 & 3.32 \\

%EGNN& 0.5497& 0.6851& 0.5519& 0.6102& 0.5992& 4.99& 2.50& 4.71& 2.88& \\
SEGNN~\citep{segnn} & 0.5327 & 0.6921 & 0.5369 & 0.6790 & 0.6101 & \textbf{5.37} & \textbf{2.46} & \textbf{4.91} & 2.63 & \textbf{3.84} \\
\midrule[0.6pt]
%Equiformer& 0.5093& 0.6285 & 0.5053 & 0.5556 & 0.5497 & 5.04 & 2.85 & 4.90 & 3.04 & \\
Equiformer & \textbf{0.5037} & \textbf{0.6881} & \textbf{0.5213} & \textbf{0.6301} & \textbf{0.5858} & 5.14 & 2.41 & 4.67 & \textbf{2.69} & 3.73 \\
%\midrule[0.6pt]
%Equiformer + NAT & 0.4737 & 0.7245 & 0.4862 & 0.6493 & 0.5834 & 6.01 & 2.45 & 5.41 & 2.71 & \\
\bottomrule[1.2pt]

\end{tabular}
}
\vspace{-3mm}
\end{adjustwidth}
\caption{\textbf{Results on OC20 IS2RE \underline{testing} set.}}
\vspace{-3mm}
\label{tab:is2re_test}
\end{table}


\begin{table}[t!]
\begin{adjustwidth}{0cm}{0cm}
\centering
\resizebox{\ifdim\width>\textwidth\textwidth\else\width\fi}{!}{
\begin{tabular}{lccccc|ccccc}
%\cline{7-11}
\toprule[1.2pt]
 & \multicolumn{5}{c|}{Energy MAE (eV) $\downarrow$} & \multicolumn{5}{c}{EwT (\%) $\uparrow$} \\ 
\cmidrule[0.6pt]{2-11}
Methods & ID & OOD Ads & OOD Cat & OOD Both & Average & ID & OOD Ads & OOD Cat & OOD Both & Average \\
\midrule[1.2pt]
GNS~\citep{noisy_nodes} & 0.54 & 0.65 & 0.55 & 0.59 & 0.5825 & - & - & - & - & -\\
%GNS + Noisy Nodes~\cite{noisy_nodes} & 0.48 & 0.53 & 0.49 & 0.48 & 0.4950 & - & - & - & - & - \\

GNS + Noisy Nodes~\citep{noisy_nodes} & 0.47 & 0.51 & 0.48 & 0.46 & 0.4800 & - & - & - & - & - \\
Graphormer~\citep{graphormer_3d} & 0.4329 & 0.5850 & 0.4441 & 0.5299 & 0.4980 & - & - & - & - & -\\
\midrule[0.6pt]
%Equiformer& 0.5093& 0.6285 & 0.5053 & 0.5556 & 0.5497 & 5.04 & 2.85 & 4.90 & 3.04 & \\
Equiformer & 0.4222 & 0.5420 & 0.4231 & 0.4754 & 0.4657 & 7.23 & 3.77 & 7.13 & 4.10 & 5.56 \\
Equiformer + Noisy Nodes & \textbf{0.4156} &  \textbf{0.4976} & \textbf{0.4165} & \textbf{0.4344} & \textbf{0.4410} & \textbf{7.47} & \textbf{4.64} & \textbf{7.19} & \textbf{4.84} & \textbf{6.04} \\
%\midrule[0.6pt]
%Equiformer + NAT & 0.4685 & 0.6071 & 0.4746 & 0.5345 & 0.5212 & 5.83 & 3.18 & 5.88 & 3.43 & \\
\bottomrule[1.2pt]
\end{tabular}
}
\end{adjustwidth}
\vspace{-3mm}
%%%\caption{\textbf{Results on OC20 IS2RE \underline{validation} set when IS2RS node-level auxiliary task is adopted during training.} 
\caption{\textbf{Results on OC20 IS2RE \underline{validation} set when IS2RS is adopted during training.} 
%%%``GNS'' denotes the 50-layer GNS trained without Noisy Nodes data augmentation, and ``GNS + Noisy Nodes'' denotes the 100-layer GNS trained with Noisy Nodes.
%Compared to the main text, we add the result of ``Equiformer + Noisy Nodes'', which use data augmentation of interpolating between initial structure and relaxed struture and adding Gaussian noise as described by Noisy Nodes~\cite{noisy_nodes}.
%%%``Equiformer + Noisy Nodes'' uses data augmentation of interpolating between initial structure and relaxed structure and adding Gaussian noise as described by Noisy Nodes~\citep{noisy_nodes}.
}
\vspace{-3mm}
\label{tab:is2re_is2rs_validation}
\end{table}


\begin{table}[t!]
\begin{adjustwidth}{0cm}{0cm}
\centering
\resizebox{\ifdim\width>\textwidth\textwidth\else\width\fi}{!}{
\begin{tabular}{lccccc|cccccb{0.15cm}b{0.15cm}}
%\cline{7-11}
\toprule[1.2pt]
 & \multicolumn{5}{c|}{Energy MAE (eV) $\downarrow$} & \multicolumn{5}{c}{EwT (\%) $\uparrow$} & \multicolumn{2}{c}{Training time} \\ 
\cmidrule[0.6pt]{2-11}
Methods & ID & OOD Ads & OOD Cat & OOD Both & Average & ID & OOD Ads & OOD Cat & OOD Both & Average & \multicolumn{2}{c}{(GPU-days)} \\
\midrule[1.2pt]
GNS + Noisy Nodes~\citep{noisy_nodes} & 0.4219 & 0.5678 & 0.4366 & \textbf{0.4651} & 0.4728 & \textbf{9.12} & \textbf{4.25} & 8.01 & \textbf{4.64} & \textbf{6.5} & \multicolumn{1}{b{0.17cm}}{56} & \multicolumn{1}{l}{(TPU)} \\
%GNS + Noisy Nodes~\cite{noisy_nodes} & 0.48 & 0.53 & 0.49 & 0.48 & 0.4950 & - & - & - & - & - \\
%
%Noisy Nodes~\cite{noisy_nodes} & 0.47 & 0.51 & 0.48 & 0.46 & 0.4800 & - & - & - & - & - \\
Graphormer~\citep{graphormer_3d}$^\dagger$ & \textbf{0.3976} & 0.5719 & \textbf{0.4166} & 0.5029 & 0.4722 & 8.97 & 3.45 & \textbf{8.18} & 3.79 & 6.1 & \multicolumn{1}{b{0.17cm}}{372} & \multicolumn{1}{l}{(A100)} \\ %372 (A100)\\
\midrule[0.6pt]
%Equiformer& 0.5093& 0.6285 & 0.5053 & 0.5556 & 0.5497 & 5.04 & 2.85 & 4.90 & 3.04 & \\
Equiformer + Noisy Nodes & 0.4171 & \textbf{0.5479} & 0.4248 & 0.4741 & \textbf{0.4660} & 7.71& 3.70 & 7.15 & 4.07 & 5.66 & \multicolumn{1}{b{0.17cm}}{\textbf{24}} & \multicolumn{1}{l}{(A6000)} \\ %\multicolumn{1}{r}{24   (A6000)} \\
%+ Noisy Nodes & 0.4156 &  0.4976 & 0.4165 & 0.4344 & 0.4410 & 7.47 & 4.64 & 7.19 & 4.84 & \\
%\midrule[0.6pt]
%Equiformer + NAT & 0.4685 & 0.6071 & 0.4746 & 0.5345 & 0.5212 & 5.83 & 3.18 & 5.88 & 3.43 & \\
\bottomrule[1.2pt]
\end{tabular}
}
\end{adjustwidth}
\vspace{-3mm}
\caption{
%\revision{
%%%\textbf{Results on OC20 IS2RE \underline{testing} set when IS2RS node-level auxiliary task is adopted during training.} 
\textbf{Results on OC20 IS2RE \underline{testing} set when IS2RS is adopted during training.} 
$\dagger$ denotes using ensemble of models trained on both IS2RE training and validation sets.
In contrast, we use the same single Equiformer model in Table~\ref{tab:is2re_is2rs_validation}, which is trained only on the training set.
Note that Equiformer achieves better results with much less computation.
%}
}
\vspace{-6mm}
\label{tab:is2re_is2rs_testing}
\end{table}


\vspace{-3mm}
\subsection{OC20}
\vspace{-2mm}
\label{subsec:oc20}

\vspace{-1mm}
\paragraph{Dataset.}
%The Open Catalyst 2020 (OC20) dataset~\cite{oc20} contains 660k distinct large atomisitc systems, each consisting of a small molecule called adsorbate placed on a large slab called catalyst.
The Open Catalyst 2020 (OC20) dataset~\citep{oc20} (Creative Commons Attribution 4.0 License) consists of larger atomic systems, each composed of a molecule called adsorbate placed on a slab called catalyst.
Each input contains more atoms and more diverse atom types than QM9 and MD17.
%%%The average number of atoms in a system is more than 70, and there are over 50 atom species.
%%%The goal is to understand interaction between adsorbates and catalysts through relaxation.
%Specifically, an adsorbate is first placed on top of a catalyst at heuristically determined locations in order to form initial structure (IS).
%%%An adsorbate is first placed on top of a catalyst to form initial structure (IS).
%%%The positions of atoms are updated with forces calculated by density function theory until the system is stable and becomes relaxed structure (RS).
%%%The energy of RS, or relaxed energy (RE), is correlated with catalyst activity and therefore a metric for understanding their interaction.
%%%We focus on the task of initial structure to relaxed energy (IS2RE), which predicts relaxed energy (RE) given an initial structure (IS).
We focus on the task of initial structure to relaxed energy (IS2RE), which is to predict the energy of a relaxed structure (RS) given its initial structure (IS).
%%%There are 460k, 100k and 100k structures in training, validation, and testing sets, respectively.
Performance is measured in MAE and energy within threshold (EwT), the percentage in which predicted energy is within 0.02~eV of ground truth energy.
In validation and testing sets, there are four sub-splits containing in-distribution adsorbates and catalysts (ID), out-of-distribution adsorbates (OOD-Ads), out-of-distribution catalysts (OOD-Cat), and out-of-distribution adsorbates and catalysts (OOD-Both). 
%We minimize MAE between predictions and normalized ground truth.
Please refer to Sec.~\ref{appendix:subsec:oc20_dataset_details} for the detailed description of OC20 dataset.

\vspace{-3.25mm}
\paragraph{Setting.}
We consider two training settings based on whether a node-level auxiliary task~\citep{noisy_nodes} is adopted.
In the first setting, we minimize MAE between predicted energy and ground truth energy without any node-level auxiliary task.
%%%In the second setting, we incorporate the task of initial structure to relaxed structure (IS2RS) as a node-level auxiliary task. %, graphormer_3d}.
%%%In addition to predicting energy, we predict node-wise vectors indicating how each atom moves from initial structure to relaxed structure.
%%%In the second setting, we incorporate the task of initial structure to relaxed structure (IS2RS) as a node-level auxiliary task and additionally predict node-wise vectors indicating how each atom moves from initial structure to relaxed structure.
In the second setting, we incorporate the task of initial structure to relaxed structure (IS2RS) as a node-level auxiliary task.

\vspace{-3.25mm}
\paragraph{Training Details.}
Please refer to Sec.~\ref{appendix:subsec:oc20_training_details} in appendix for details on Equiformer architecture, hyper-parameters and training time.


\vspace{-3.25mm}
\paragraph{IS2RE Results without Node-Level Auxiliary Task.}
We summarize the results on validation and testing sets under the first setting in Table~\ref{tab:is2re_validation} in appendix and Table~\ref{tab:is2re_test}.
Compared with state-of-the-art models like SEGNN and SphereNet, Equiformer consistently achieves the lowest MAE for all the four sub-splits in validation and testing sets.
%%%Note that energy within threshold (EwT) considers only the percentage of predictions close enough to ground truth and the distribution of errors, and therefore improvement in average errors (MAE) would not necessarily reflect that in error distributions (EwT).
Note that EwT considers only the percentage of predictions close enough to ground truth and the distribution of errors, and therefore improvement in average errors (MAE) would not necessarily reflect that in error distributions (EwT).
A more detailed discussion can be found in Sec.~\ref{appendix:subsec:error_distributions} in appendix.
%as long as errors are still higher than 0.02~eV.
%Similar phenomena can be observed in OC20 leaderboard\footnote{As shown in the leaderboard (\href{https://opencatalystproject.org/leaderboard.html}{https://opencatalystproject.org/leaderboard.html}), ``GNS + Noisy Nodes (Direct)'' achieves lower MAE but lower EwT than ``DimeNet++-1.8M-fonly-eonly-relaxation-All''. 
%A screenshot of the leaderboard can be found \href{https://drive.google.com/file/d/1nxS_Zb4VrV1nIEAhkPFaj9mcqaHfQiL_/view?usp=sharing}{here}.}.
%%%Similar phenomena can be observed in Table~\ref{tab:is2re_test}, where for ``OOD Both'' sub-split, SphereNet achieves lower MAE yet lower EwT than SEGNN.
%%%We also note that models in Table~\ref{tab:is2re_validation} and~\ref{tab:is2re_test} are trained by minimizing MAE, and therefore comparing MAE in validation and testing sets could mitigate the discrepancy between training objectives and evaluation metrics and that OC20 leaderboard ranks the relative performance of models mainly according to MAE.
We also note that models are trained by minimizing MAE, and therefore comparing MAE could mitigate the discrepancy between training objectives and evaluation metrics and that OC20 leaderboard ranks the relative performance according to MAE.
Additionally, we compare the training time of SEGNN and Equiformer in Sec.~\ref{appendix:subsec:oc20_training_time_comparison}.

\vspace{-3.5mm}
\paragraph{IS2RE Results with IS2RS Node-Level Auxiliary Task.}
We report the results on validation and testing sets in Table~\ref{tab:is2re_is2rs_validation} and Table~\ref{tab:is2re_is2rs_testing}.
As of the date of the submission of this work, Equiformer achieves the best results on IS2RE task when only IS2RE and IS2RS data are used.
%%%Notably, the proposed Equiformer in Table~\ref{tab:is2re_is2rs_testing} achieves competitive results even with much less computation.
Notably, the result in Table~\ref{tab:is2re_is2rs_testing} is achieved with much less computation.
%Specifically, training ``Equiformer + Noisy Nodes'' takes about $24$ GPU-days when A6000 GPUs are used.
%The training time of ``GNS + Noisy Nodes'' is $56$ TPU-days.
%``Graphormer'' uses ensemble of $31$ models and requires $372$ GPU-days when A100 GPUs are used.
%The comparison to GNS demonstrates the improvement from invariant message passing networks to equivariant Transformers.
%Without any data augmentation, Equiformer still achieves competitive results to GNS trained with Noisy Nodes~\cite{noisy_nodes}.
%Compared to Graphormer, Equiformer demonstrates the effectiveness of equivariant features and the proposed equivariant graph attention.
We note that under this setting, greater depths and thus more computation translate to better performance~\citep{noisy_nodes} and that Equiformer demonstrates incorporating equivariant features and the proposed equivariant graph attention can improve training efficiency by $2.3\times$ to $15.5\times$ compared to invariant message passing networks and invariant Transformers.
%%%Note that Equiformer, with $18$ Transformer blocks, is relatively shallow as GNS trained with Noisy Nodes has $100$ blocks and Graphormer has $48$ Transformer blocks and that deeper networks can typically obtain better results when IS2RS auxiliary task is adopted~\cite{noisy_nodes}.


\vspace{-2mm}

\begin{table}[t]
\centering
\resizebox{\ifdim\width>\textwidth\textwidth\else\width\fi}{!}{

\begin{tabular}{cccclcccccccccc}
\toprule[1.2pt]
& \multicolumn{3}{c}{Methods} & & \\
\cmidrule[0.6pt]{2-4}
& \multirow{2}{*}{\begin{tabular}[c]{@{}c@{}}Non-linear \\ message passing\end{tabular}} & \multirow{2}{*}{\begin{tabular}[c]{@{}c@{}}MLP \\ attention\end{tabular}} & \multirow{2}{*}{\begin{tabular}[c]{@{}c@{}}Dot product \\ attention\end{tabular}} &  Task & $\alpha$ & $\Delta \varepsilon$ & $\varepsilon_{\text{HOMO}}$ & $\varepsilon_{\text{LUMO}}$ & $\mu$ & $C_{\nu}$ & & Training time & Number of \\
%Index & & & & Unit & bohr$^3$ & meV & meV & meV & D & cal/mol K & & (minutes/epoch) & parameters \\
Index & & & & Unit & $a_0^3$ & meV & meV & meV & D & cal/mol K & & (minutes/epoch) & parameters \\
\midrule[1.2pt]

1 & \ding{51} & \ding{51} &  & & .046 & 30 & 15 & 14 & .011 & .023 & & 12.1 & 3.53M \\
2 & & \ding{51} &   & & .051 & 32 & 16 & 16 & .013 & .025 & & 7.2 & 3.01M \\
3 & & & \ding{51} & & .053 & 32 & 17 & 16 & .013 & .025 & & 7.8 & 3.35M \\
%4 & & \ding{51} &  &  & & .060 &  & 18 & 17 & .016 &  \\

\bottomrule[1.2pt]

\end{tabular}
}
\vspace{-3mm}
\caption{\textbf{Ablation study results on QM9.}}
\vspace{-3.25mm}
\label{tab:qm9_ablation}
\end{table}


\begin{table}[t]
\begin{adjustwidth}{0cm}{0cm}
\centering
\resizebox{\ifdim\width>\textwidth\textwidth\else\width\fi}{!}{

\begin{tabular}{cccccccccccccccccc}
\toprule[1.2pt]
& \multicolumn{3}{c}{Methods} & & \multicolumn{5}{c}{Energy MAE (eV) $\downarrow$} & & \multicolumn{5}{c}{EwT (\%) $\uparrow$} \\
\cmidrule[0.6pt]{2-4}
\cmidrule[0.6pt]{6-10}
\cmidrule[0.6pt]{12-16}
& \multirow{2}{*}{\begin{tabular}[c]{@{}c@{}}Non-linear \\ message passing\end{tabular}} & \multirow{2}{*}{\begin{tabular}[c]{@{}c@{}}MLP \\ attention\end{tabular}} & \multirow{2}{*}{\begin{tabular}[c]{@{}c@{}}Dot product \\ attention\end{tabular}} & & & & & & & & & & & & & Training time & Number of \\

Index & & & &  & ID & OOD Ads & OOD Cat & OOD Both & Average & & ID & OOD Ads & OOD Cat & OOD Both & Average & (minutes/epoch) & parameters \\
\midrule[1.2pt]

1 & \ding{51} & \ding{51} &  & & 0.5088 & 0.6271 & 0.5051 & 0.5545 &  0.5489  & & 4.88 & 2.93 & 4.92 & 2.98 & 3.93 & 130.8 & 9.12M \\
2 & & \ding{51} &   & & 0.5168 & 0.6308 & 0.5088 & 0.5657 & 0.5555 &  & 4.59 & 2.82 & 4.79 & 3.02 & 3.81 & 91.2 & 7.84M \\
3 & & & \ding{51}  & & 0.5386 & 0.6382 & 0.5297 & 0.5692 & 0.5689 & & 4.37 & 2.60 & 4.36 & 2.86 & 3.55 & 99.3 & 8.72M \\
%4 & & \ding{51} &  &  & 0.5213 & 0.6268 & 0.5178 & 0.5621 & 0.5579\\

%\bottomrule[1.2pt]
%\bottomrule[1.2pt]
%\vspace{-1mm}
%\\
%\toprule[1.2pt]

%%%\midrule[1.2pt]
%%%& \multicolumn{3}{c}{Methods} & & \multicolumn{5}{c}{EwT (\%) $\uparrow$} \\
%%%\cmidrule[0.6pt]{2-4}
%%%\cmidrule[0.6pt]{6-10}
%%%& \multirow{2}{*}{\begin{tabular}[c]{@{}c@{}}Non-linear \\ message passing\end{tabular}} & \multirow{2}{*}{\begin{tabular}[c]{@{}c@{}}MLP \\ attention\end{tabular}} & \multirow{2}{*}{\begin{tabular}[c]{@{}c@{}}Dot product \\ attention\end{tabular}} &  \\

%%%Index & & & & & ID & OOD Ads & OOD Cat & OOD Both & Average \\
%%%\midrule[1.2pt]

%%%1 & \ding{51} & \ding{51} & & & 4.88 & 2.93 & 4.92 & 2.98 & 3.93 \\
%%%2 & & \ding{51} &   &  & 4.59 & 2.82 & 4.79 & 3.02 & 3.81  \\
%%%3 & & & \ding{51} & & 4.37 & 2.60 & 4.36 & 2.86 & 3.55 \\
%4 & & \ding{51} &  &  & 4.37 & 2.80 & 4.40 & 2.98 & 3.64 \\

\bottomrule[1.2pt]

\end{tabular}
}

\vspace{-3mm}
\end{adjustwidth}
\caption{\textbf{Ablation study results on OC20 IS2RE validation set.}
%\note{Dot product attention uses smaller learning rate $1.5 \times 10^{-3}$.}
}
\vspace{-6.75mm}
\label{tab:oc20_ablation}
\end{table}



\vspace{-2.25mm}
\subsection{Ablation Study}
\vspace{-2.5mm} 
\label{subsec:ablation_study}

We conduct ablation studies \revision{to show that Equiformer with dot product attention and linear message passing has already achieved strong empirical results and demonstrate} the improvement brought by MLP attention and non-linear messages in the proposed equivariant graph attention.
%We modify dot product (DP) attention~\citep{transformer, se3_transformer} so that it only differs from MLP attention in how attention weights $a_{ij}$ are generated from $f_{ij}$.
\revision{Dot product (DP) attention only differs from MLP attention in how attention weights $a_{ij}$ are generated from $f_{ij}$.}
%%%Please refer to Sec.~\ref{appendix:subsec:dot_product_attention} in appendix for details on DP attention.
Please refer to Sec.~\ref{appendix:subsec:dot_product_attention} in appendix for further details.
For experiments on QM9 and OC20, unless otherwise stated, we follow the hyper-parameters used in previous experiments. 
%\revision{Please refer to Sec.~\ref{appendix:subsec:qm9_training_details} and Sec.~\ref{appendix:subsec:oc20_training_details} for training time of different models.}

\vspace{-3.25mm}
\paragraph{Result on QM9.}
The comparison is summarized in Table~\ref{tab:qm9_ablation}.
\revision{Compared with models in Table~\ref{tab:qm9_result}, Equiformer with dot product attention and linear message passing (Index 3) achieves competitve results.}
%Non-linear messages improve upon linear messages when MLP attention is used.
Non-linear messages improve upon linear messages when MLP attention is used while non-linear messages increase the number of tensor product operations in each block from $1$ to $2$ and thus inevitably increase training time.
%%%Similar to what is reported by GATv2~\cite{gatv2}, the improvement of replacing DP attention with MLP attention is not very significant.
%%%\todo{On the other hand, the improvement of replacing DP attention with MLP attention is not very significant.}
On the other hand, MLP attention achieves similar results to DP  attention.
We conjecture that DP attention with linear operations is expressive enough to capture common attention patterns as the numbers of nighboring nodes and atom species are much smaller than those in OC20.
However, MLP attention is roughly $8\%$ faster as it directly generates scalar features and attention weights from $f_{ij}$ instead of producing additional key and query irreps features for attention weights.

\vspace{-3.25mm}
\paragraph{Result on OC20.}
%%%We consider the setting of training without IS2RS auxiliary task and use a smaller learning rate $1.5 \times 10 ^{-4}$ for DP attention as this improves the performance.
We consider the setting of training without auxiliary task and summarize the comparison in Table~\ref{tab:oc20_ablation}.
%%%We summarize the comparison in Table~\ref{tab:oc20_ablation}.
\revision{Compared with models in Table~\ref{tab:is2re_validation}, Equiformer with dot product attention and linear message passing (Index 3) has already outperformed all previous models.}
Non-linear messages consistently improve upon linear messages.
In contrast to the results on QM9, MLP attention achieves better performance than DP attention and is $8\%$ faster.
We surmise this is because OC20 contains larger atomistic graphs with more diverse atom species and therefore requires more expressive attention mechanisms.
Note that Equiformer can potentially improve upon previous equivariant Transformers~\citep{se3_transformer, torchmd_net, eqgat} since they use less expressive attention mechanisms similar to Index 3 in Table~\ref{tab:oc20_ablation}.